% Preámbulo del informe de la Etapa 2.
%
% Tipografía: Computer Modern a 12pt, que es exactamente la que usa la cátedra
% en «Reglas de Sintaxis de MiniJava» (pdffonts sobre ese PDF reporta CMR12 para
% el cuerpo, CMCSC10 para las versalitas, CMTT10 para las producciones y
% CMSSBX10 para los terminales). `lmodern` es Latin Modern, la reimplementación
% Unicode del mismo diseño.
%
% Paleta monocromática y estructura heredadas del informe de la Etapa 1.
\usepackage[T1]{fontenc}
\usepackage[utf8]{inputenc}
\usepackage{lmodern}
\usepackage{amsmath,amssymb}
\usepackage[spanish,es-noquoting,es-noshorthands,es-tabla]{babel}
\usepackage[a4paper,margin=2.6cm,headsep=14pt]{geometry}
\usepackage{microtype}
\usepackage{booktabs}
\usepackage{longtable}
\usepackage{array}
\usepackage{xcolor}
\usepackage{fancyhdr}
\usepackage{titlesec}
\usepackage{enumitem}
\usepackage{fancyvrb}
\usepackage{framed}
\usepackage[hidelinks,bookmarks=true]{hyperref}

% Caracteres que aparecen en los comentarios de los archivos de gramática y
% llegan al informe a través de los listados generados.
\DeclareUnicodeCharacter{03B5}{\ensuremath{\epsilon}}
\DeclareUnicodeCharacter{2014}{---}
\DeclareUnicodeCharacter{2026}{\ldots}

% --- Paleta monocromática: sólo grises, con un único acento. -----------------
\definecolor{tinta}{gray}{0.13}
\definecolor{grisMedio}{gray}{0.42}
\definecolor{grisSuave}{gray}{0.70}
\definecolor{grisFondo}{gray}{0.955}
\definecolor{acento}{HTML}{2F4858}

\color{tinta}

% --- Títulos ----------------------------------------------------------------
\titleformat{\section}{\Large\bfseries\color{acento}}{\thesection}{0.7em}{}
\titleformat{\subsection}{\large\bfseries\color{tinta}}{\thesubsection}{0.7em}{}
\titleformat{\subsubsection}{\bfseries\color{grisMedio}}{\thesubsubsection}{0.7em}{}
\titlespacing*{\section}{0pt}{2.4ex plus 1ex minus .2ex}{1.3ex plus .2ex}

% --- Encabezado y pie -------------------------------------------------------
\setlength{\headheight}{14pt}
\pagestyle{fancy}
\fancyhf{}
\renewcommand{\headrulewidth}{0.4pt}
\renewcommand{\headrule}{\hbox to\headwidth{\color{grisSuave}\leaders\hrule height \headrulewidth\hfill}}
\fancyhead[L]{\footnotesize\color{grisMedio}Compiladores e Intérpretes — Etapa 2}
\fancyhead[R]{\footnotesize\color{grisMedio}\nouppercase{\leftmark}}
\fancyfoot[C]{\footnotesize\color{grisMedio}\thepage}

% ===========================================================================
% Notación de la gramática — la misma convención tipográfica de la cátedra
% ===========================================================================
% \nt{X}   no terminal   -> \texttt, como <X> en el PDF de la cátedra
% \tl{x}   terminal      -> sans-serif negrita, como los terminales del PDF
% \eps     la cadena vacía
% \deriva  el ::= de las producciones
\newcommand{\nt}[1]{\texttt{<#1>}}
\newcommand{\tl}[1]{\textsf{\bfseries #1}}
\newcommand{\eps}{\ensuremath{\epsilon}}
\newcommand{\deriva}{\texttt{::=}}
\newcommand{\alt}{\ensuremath{\;\mid\;}}
\newcommand{\prim}[1]{\textup{Primeros}(#1)}

% Bloque de producciones: sangrado, interlineado propio, sin justificar.
\newenvironment{producciones}{%
  \par\smallskip\begingroup
  \parindent=0pt \leftskip=1.2em
  % Elasticidad acotada: evita estirar los espacios de una producción larga,
  % pero no tanta como para que el \hfill de la etiqueta de logro no empuje.
  \rightskip=0pt plus 2em
  \parfillskip=0pt plus 1fil
  \setlength{\parskip}{0.35ex}%
  \obeylines}{\par\endgroup\smallskip}

% Una producción con su nombre a la izquierda, para las tablas de Primeros.
\newcommand{\regla}[2]{\nt{#1} \deriva{} #2}

% --- Bloques de código y de consola -----------------------------------------
\DefineVerbatimEnvironment{consola}{Verbatim}%
  {fontsize=\small,xleftmargin=1.1em,formatcom=\color{tinta}}
\DefineVerbatimEnvironment{codigo}{Verbatim}%
  {fontsize=\small,xleftmargin=1.1em,commandchars=\\\{\}}

% --- Cajas de defecto -------------------------------------------------------
% Dos clases de defecto que NO hay que confundir; van con marcas distintas.
\newenvironment{conflicto}[1]{%
  \par\smallskip\begingroup
  \setlength{\fboxsep}{0.7em}%
  \definecolor{shadecolor}{gray}{0.955}%
  \begin{shaded}\noindent{\bfseries\color{acento}Conflicto LL(1) --- #1}\par\smallskip
}{\end{shaded}\endgroup\smallskip}

\newenvironment{ambiguedad}[1]{%
  \par\smallskip\begingroup
  \setlength{\fboxsep}{0.7em}%
  \definecolor{shadecolor}{gray}{0.90}%
  \begin{shaded}\noindent{\bfseries\color{tinta}Ambigüedad --- #1}\par\smallskip
}{\end{shaded}\endgroup\smallskip}

\newenvironment{decision}[1]{%
  \par\smallskip\begingroup
  \setlength{\fboxsep}{0.7em}%
  \definecolor{shadecolor}{gray}{0.975}%
  \begin{shaded}\noindent{\bfseries\color{grisMedio}Decisión --- #1}\par\smallskip
}{\end{shaded}\endgroup\smallskip}

% --- Los cuatro pasos de cada transformación --------------------------------
\newcommand{\paso}[1]{\par\smallskip\noindent\textbf{\color{grisMedio}#1}\quad\ignorespaces}

% Los nombres largos de no terminales van en \texttt y no se parten, así que se
% le da al justificado un margen de maniobra en vez de dejar renglones desbordados.
\setlength{\emergencystretch}{2.5em}

\setlist[itemize]{leftmargin=1.3em,itemsep=0.25ex,topsep=0.5ex}
\setlist[enumerate]{leftmargin=1.5em,itemsep=0.25ex,topsep=0.5ex}
\setlist[description]{leftmargin=1.3em,itemsep=0.35ex,topsep=0.6ex}
\renewcommand{\arraystretch}{1.25}
